\documentclass[conference]{IEEEtran}

\usepackage{cite}
\usepackage{graphicx}
\usepackage{booktabs}
\usepackage{url}
\usepackage{amsmath}
\usepackage[hidelinks]{hyperref}

\renewcommand{\figurename}{Fig.}
\graphicspath{{../figures/}}

\begin{document}

\title{Do Systematic Literature Reviews Cite the Right Papers? \\
       A Citation-Coverage Audit of Technical-Debt SLRs%
       \thanks{Source code, classifications, and intermediate data are
       available at \url{https://github.com/connortessaro/slr-citation-audit}.}}

\author{%
  \IEEEauthorblockN{Connor Tessaro, Lloyd Thomas, Jaden Zhou}
  \IEEEauthorblockA{Khoury College of Computer Sciences \\
                    Northeastern University \\
                    Boston, MA, USA \\
                    \texttt{\{tessaro.c, thomas.ll, zhou.jad\}@northeastern.edu}}
}

\maketitle

\begin{abstract}
Systematic Literature Reviews (SLRs) claim to deliver comprehensive, unbiased
coverage of a research area. This study asks a narrower but consequential
question: do SLRs in the technical-debt subfield actually cite the papers
that the wider community most heavily relies on? We identified 60
peer-reviewed SLRs and systematic mapping studies published between 2015 and
2026, extracted their reference lists from Semantic Scholar with a Crossref
fallback, and independently constructed a two-pass top-50 most-cited paper
list using Semantic Scholar citation counts. With year-of-publication date
controls applied, mean coverage of the top-cited corpus across SLRs was
10.5\,\% (median 4.4\,\%, range 0--40\,\%); 24 of 60 SLRs (40\,\%) cited
none of the eligible top-50. Gaps clustered in journal venues (90.4\,\% of
1{,}859 missed pairs) and showed a dual age peak---papers 6--10 years old at
SLR time accounted for 34.8\,\% of misses, and papers 0--2 years old for a
further 32.0\,\%. Forty-three percent of missed papers were open-access,
suggesting access barriers explain only a limited share of the gap. The
headline finding is robust under a top-100 benchmark (mean 8.6\,\%). The
analysis is Semantic Scholar end-to-end; ACM exports were collected for
traceability but not used in the reported overlap. We discuss methodological
implications and offer protocol-level recommendations to improve coverage
in future SLRs.
\end{abstract}

\begin{IEEEkeywords}
systematic literature review, bibliometrics, technical debt, citation
coverage, research methodology
\end{IEEEkeywords}

\section{Introduction}

SLRs are positioned as the most rigorous form of literature synthesis in
software engineering, following protocols such as Kitchenham \&
Charters~\cite{kitchenham2007guidelines} and the PRISMA reporting
guidelines~\cite{moher2009prisma}. They promise to surface a representative,
unbiased view of a research area. If that promise holds, the most
consequential primary studies---those the community most heavily
cites---should appear in the reference lists of well-conducted SLRs.

This study tests that promise within a single, well-defined subfield:
technical debt. Technical debt is an attractive testbed because it has a
strong SLR tradition since 2015 (with the metaphor itself dating to
Cunningham~\cite{cunningham1992wycash}), it spans more than a decade of
primary literature, and its scope is reasonably bounded by a small set of
multi-word phrases.

During the project we revised the identification methodology three
times---dropping a ``methodology-signal'' gate that had been excluding
non-PRISMA-worded SLRs, adding a Crossref backfill for SLRs whose
references were not indexed by Semantic Scholar, and widening the
secondary-study self-label patterns to cover terms such as \emph{literature
review}, \emph{scoping review}, and \emph{survey} alongside the strict
\emph{systematic \ldots} family. Each revision is documented under
\texttt{docs/changes/} with rationale and effect on corpus size. We surface
them so readers can judge the inclusion criteria transparently.

Our research questions are:

\begin{itemize}
  \item \textbf{RQ1.} What proportion of the top-50 most-cited
        technical-debt papers (with date controls) appears in each published
        SLR?
  \item \textbf{RQ2.} Where coverage gaps exist, are they explained by
        venue type, publication age relative to the SLR, open-access
        status, or DOI block (ACM vs.\ IEEE)?
  \item \textbf{RQ3.} What practical implications do the observed gaps
        carry for how SLRs in this subfield should be conducted?
\end{itemize}

\section{Method}

\subsection{Subfield definition}
We treated technical debt as the set of papers using the terms
\emph{technical debt}, \emph{design debt}, \emph{code debt}, or
\emph{architectural debt} as the primary topic. Concrete inclusion criteria
are documented in \texttt{data/manual/CLASSIFICATION.md}.

\subsection{SLR identification}
We queried Semantic Scholar via 16 API searches over 2000--2026---each of
four subfield keywords crossed with each of three systematic-review phrases,
plus four broad keyword-only searches with a title-pattern post-filter. ACM
Digital Library bronze BibTeX exports were collected for traceability,
parsed, and deduplicated against the Semantic Scholar pool, but the
\emph{reported overlap is computed from Semantic Scholar alone}. ACM and
IEEE end-to-end overlap runs are noted as future cross-validation, not as
part of this analysis. All candidates were deduplicated on a key derived
from DOI, then Semantic Scholar \texttt{paperId}, then a normalised title.

\subsection{SLR classification}
A candidate is included in the SLR corpus only if all of the following hold:
(1)~the title or abstract mentions a subfield keyword; (2)~the title or
abstract contains a secondary-study self-label (we widened this pattern in
June 2026 to include the systematic family plus \emph{literature review},
\emph{scoping review/study}, \emph{mapping study}, \emph{SLR}, \emph{survey},
\emph{tertiary review/study}, and \emph{meta-analysis}); (3)~the publication
year falls in 2000--2026; and (4)~the candidate has an extractable reference
list after Semantic Scholar extraction and Crossref fallback. SLRs with no
bibliography are excluded because coverage is undefined without references.

\emph{Trade-off on widened labels.} Patterns such as \emph{survey} and
\emph{literature review} raise recall but admit narrative reviews without
strict protocols. We accept that for this audit because (i)~the
subfield-keyword gate removes most off-topic noise, and (ii)~the research
question is bibliographic coverage, not quality appraisal of review
methodology. A poorly-conducted narrative survey that omits influential
citations is still informative for RQ1. We re-state this trade-off as a
limitation in Section~\ref{sec:limitations}.

\subsection{Reference extraction}
For each included SLR we queried Semantic Scholar's
\texttt{/paper/\{id\}/references} endpoint, paginating to exhaustion. Each
cited paper was normalised to a \texttt{(paper\_key, title, year, venue,
doi, citationCount, openAccessPdf)} record, and per-SLR results were cached
locally so re-runs do not re-hit the API. Where Semantic Scholar returned
an empty reference list, we attempted a Crossref backfill against the SLR's
DOI. SLRs that still had no bibliography after both sources were excluded.
The resulting funnel---8{,}417 raw SS candidates $\to$ 110 auto-INCLUDE $\to$
60 SLRs in the analysis corpus---is reported in
Section~\ref{sec:corpus-overview}.

\subsection{Top-cited corpus construction}
The Semantic Scholar search API does not support ``rank by citations''
directly. We issued a wide keyword search ($\textit{limit}=1{,}000$ per
keyword), pooled the results, deduplicated by \texttt{paper\_key}, and
filtered to papers whose title or abstract included a subfield keyword. The
discovery pool contained 958 papers.

To avoid a single lifetime-citation ranking favouring older papers, we
constructed the benchmark with a two-pass design:

\begin{enumerate}
  \item \textbf{Established pass}---papers with publication year $\leq$
        $\textit{as\_of\_year}-4$ (i.e.~$\leq$~2022 when $\textit{as\_of\_year}=2026$),
        ranked by Semantic Scholar \texttt{citationCount}, take 25. Pool
        size 781.
  \item \textbf{Recent pass}---papers from the subsequent four calendar
        years, same ranking, take 25. Pool size 177.
\end{enumerate}

Citations across the top-50 range from 1{,}437 at rank~1 to 12 at rank~50. A
parallel two-pass list with $N=100$ (rank-100 cutoff: 9 citations) supports
the robustness check in Section~\ref{sec:coverage}.

\emph{Why not a recent-only benchmark.} Skewing the benchmark toward recent
papers would change the research question from ``do SLRs cite canonical TD
literature'' to ``do SLRs cite recent hot papers.'' The two-pass design
balances accumulated influence and recent momentum without abandoning the
older highly-cited work that defines the subfield.

\subsection{Overlap computation with date controls}
For each SLR, the top-cited corpus was filtered to papers with publication
year $\leq$ the SLR's publication year. This date control prevents counting
a paper as ``missed'' if it did not yet exist when the SLR was written. We
then computed
\begin{equation}
  \mathrm{coverage}_i = \frac{\mathrm{hits}_i}{\mathrm{eligible\_top\_n}_i}
  \label{eq:coverage}
\end{equation}
where $\mathrm{hits}_i$ is the number of date-eligible top-cited papers
appearing in SLR~$i$'s references. Reference matching uses
\texttt{paper\_key} (DOI $\to$ SS paperId $\to$ normalised title).

\subsection{Gap explanation}
For each $(\textit{SLR}, \textit{missed top-cited paper})$ pair we attached
the following features: \texttt{venue\_type} (journal / conference /
workshop / preprint / book / unknown), \texttt{year\_delta}, \texttt{age\_bucket}
(0--2y / 3--5y / 6--10y / 10y+ / post-slr / unknown), \texttt{open\_access},
\texttt{in\_acm} / \texttt{in\_ieee} (DOI prefix), and \texttt{is\_top10}.

\subsection{Tooling}
All analysis is implemented in Python and reproducible from the per-source
pipeline scripts. The Semantic Scholar client wrapper
(\texttt{core/ss\_client.py}) persistently caches API responses so the full
pipeline can be re-run idempotently. The interactive overlap explorer
(\texttt{tools/explorer/}) recomputes coverage in-browser when the benchmark
size or SLR cohort toggle changes.

\section{Results}

\subsection{SLR corpus overview}
\label{sec:corpus-overview}

We identified 8{,}417 unique Semantic Scholar candidates and 110 met the
inclusion criteria after the classification gates. After reference
extraction with Crossref fallback, 50 of those 110 candidates had no
extractable bibliography and were dropped; the final analysis corpus is
$N=60$ SLRs. The corpus spans 2015--2026 and is composed of 39 systematic
literature reviews, 19 systematic mapping studies, and 2 surveys.

\emph{Field-emergence finding.} Of the 8{,}417 SS candidates, 1{,}875 fall
in 2000--2014. \emph{None} of these pass both the subfield-keyword gate and
the (widened) secondary-study self-label gate. The pre-2015 records in the
pool are primary empirical studies, workshop short papers, and position
pieces---not self-labelled syntheses. We treat this as a finding about the
field rather than a pipeline failure: the technical-debt SLR tradition
emerges around 2015, after Li et al.~\cite{li2015mapping} established the
subfield's secondary-literature canon.

\subsection{Coverage of top-cited papers}
\label{sec:coverage}

Across 60 SLRs the date-controlled coverage of the top-50 most-cited
technical-debt papers ranged from 0\,\% to 40\,\%, with a mean of 10.5\,\%
and a median of 4.4\,\% (Fig.~\ref{fig:coverage}). The distribution is
right-skewed: most SLRs cite only a small fraction of the eligible top-50,
while a minority cluster above 20\,\%.

\begin{figure}[t]
  \centering
  \includegraphics[width=\columnwidth]{ss/coverage_histogram.png}
  \caption{Distribution of SLR coverage of the date-controlled top-50
           technical-debt corpus.}
  \label{fig:coverage}
\end{figure}

A striking subgroup shows zero coverage: 24 of 60 SLRs (40\,\%) cite none
of the eligible top-50 papers in their reference lists. This is the modal
outcome, not a tail effect.

\textbf{Cohort effect.} Coverage falls sharply with SLR publication year
(Table~\ref{tab:cohort}). The earlier cohort averages roughly 2.5$\times$
the coverage of the later one. The effect is partly mechanical (by later
publication years more benchmark papers are date-eligible, so the
denominator grows faster than any given SLR's reference list) and partly a
quality signal (the foundational 2015--2019 TD SLRs were written when the
canonical primary literature was smaller and more cohesive).

\begin{table}[t]
  \centering
  \caption{Coverage by SLR publication-year cohort}
  \label{tab:cohort}
  \begin{tabular}{lrrr}
    \toprule
    Cohort & $n$ & Mean & Median \\
    \midrule
    2015--2019 & 11 & 20.3\,\% & 25.0\,\% \\
    2020--2026 & 49 & 8.3\,\% & 2.0\,\% \\
    \bottomrule
  \end{tabular}
\end{table}

\textbf{Robustness.} A parallel run with the top-100 benchmark produces a
mean coverage of 8.6\,\% (median 3.5\,\%); the headline finding is robust
to corpus size. Mean coverage drops because the larger denominator grows
faster than additional hits.

\subsection{Gap characterisation}
\label{sec:gaps}

We enriched 1{,}859 $(\textit{SLR}, \textit{missed paper})$ pairs with
venue, age, and access features (Fig.~\ref{fig:gap-heatmap}).

\begin{figure}[t]
  \centering
  \includegraphics[width=\columnwidth]{ss/gap_heatmap.png}
  \caption{Missed highly-cited papers by venue type and age bucket at SLR
           publication time.}
  \label{fig:gap-heatmap}
\end{figure}

The most-missed venue type is \textbf{journal}, accounting for 90.4\,\%
(1{,}681 of 1{,}859) of all gaps; preprints are 3.1\,\% (58), and an
unknown bucket accounts for the remaining 6.5\,\% (120). Conference and
workshop misses are negligible after venue-type normalisation. The
journal-dominance result argues against a simple ``conference-visibility''
explanation for missed citations.

\textbf{Dual age peak.} Missed papers cluster at two ages relative to SLR
publication: 6--10 years old at SLR time accounts for 34.8\,\% (647
missed pairs); 0--2 years old accounts for 32.0\,\% (594); 3--5 years for
18.5\,\% (343); 10+ years for 14.8\,\% (275). The 6--10y peak is consistent
with an ``assumed-known foundations'' pattern (older canonical work treated
as background rather than cited explicitly). The 0--2y peak is consistent
with a ``recency penalty'' (newly published work missed by the SLR's
search snapshot or not yet visible enough). Both stories appear in the
same data.

\textbf{Open-access status.} 43.0\,\% (800) of missed papers had an
\texttt{openAccessPdf.url} recorded by Semantic Scholar. Access barriers
therefore explain only a limited share of the gap: SLRs are missing
freely-readable canonical papers nearly as often as paywalled ones.

\textbf{Top-10 subset.} Of the 1{,}859 missed pairs, 484 involve papers in
the global top-10 of the citation ranking---the most influential primary
studies are being missed at scale, not merely the long tail.

\section{Discussion}

\subsection{Plausible drivers of the gap}

Several mechanisms could plausibly produce the gap we observe:

\begin{itemize}
  \item \textbf{Venue invisibility.} If SLRs systematically draw from a
        limited set of digital libraries, papers in venues outside those
        libraries will be missed even when they are highly cited
        downstream. Our 90.4\,\%-journal finding constrains this story: the
        misses are concentrated in mainstream peer-reviewed venues, not in
        obscure proceedings.
  \item \textbf{Recency penalty and assumed-known foundations.} The dual
        age peak suggests two mechanisms operating together. Papers
        published shortly before an SLR may not have accumulated enough
        citations or visibility to be picked up by the SLR's search
        protocol; conversely, older foundational papers may be implicitly
        treated as background and omitted. Neither effect alone explains
        the pattern.
  \item \textbf{Protocol blind spots.} Even well-documented protocols
        depend on keyword choices that may miss papers using adjacent
        terminology (e.g.\ \emph{design debt} vs.\ \emph{technical debt})
        or that pre-date the dominant subfield label.
\end{itemize}

The cohort effect in Section~\ref{sec:coverage} is partly mechanical and
partly a quality signal. We treat it as evidence that newer SLRs are not
necessarily worse-conducted, but that the bar for ``comprehensive coverage''
has risen with the field. The journal-dominance, 43\,\%-OA, and top-10-miss
findings together suggest the gap is \emph{not} primarily explained by
venue access, paywalls, or obscurity of the missed papers. It is a
citation-protocol gap, not an information-access gap.

\subsection{What this is not}

We cannot determine from this analysis whether the \emph{conclusions} of any
individual SLR are wrong. A gap in citation coverage does not imply a
flawed synthesis: an SLR may legitimately exclude a paper that is
highly-cited but methodologically off-topic for its scope. The finding is
at the level of \emph{input coverage}, not \emph{output validity}.
Subsequent work would need to read each missed top-cited paper against
each SLR's inclusion criteria to determine whether the omission was
substantively justified.

\section{Implications for SLR practice}

Four concrete recommendations follow, each grounded in a specific
evidentiary point above:

\begin{enumerate}
  \item \textbf{Backward snowballing is non-optional.} Adding at least one
        round of backward snowballing from a small seed set of well-known
        papers catches venue-invisible work that keyword search
        misses~\cite{wohlin2014snowballing}. The field-emergence finding
        (Section~\ref{sec:corpus-overview}) supports this directly:
        pre-2015 technical-debt material exists in the candidate pool but
        does not surface via keyword + secondary-study search.
  \item \textbf{Citation-aware quality check.} After applying inclusion
        criteria, cross-check the candidate set against a top-$N$
        most-cited list in the subfield; flag any high-citation paper that
        the protocol excludes and document the reason. The 24/60
        zero-coverage finding shows that 40\,\% of SLRs would have flagged
        a non-empty set of missing canonical work.
  \item \textbf{Multi-source dedup.} Use at least three digital libraries
        plus a citation-graph source (e.g.\ Semantic Scholar) to mitigate
        single-source coverage gaps. Our SS-only scope makes this
        recommendation especially relevant---we cannot claim our headline
        number is what would have emerged from a fully cross-validated
        multi-source run.
  \item \textbf{Time-window honesty.} Distinguish ``we excluded this
        paper'' from ``this paper post-dated our search.'' The cohort
        effect shows that newer SLRs face a structurally larger
        eligibility set; readers cannot assess validity without that
        distinction made explicit.
\end{enumerate}

\section{Limitations}
\label{sec:limitations}

\begin{itemize}
  \item \textbf{Single citation source.} We use Semantic Scholar for
        citation counts and reference extraction. ACM, IEEE, and Google
        Scholar produce different counts; our top-50 may differ from a
        top-50 built from another source. ACM bronze exports were
        collected but not analysed end-to-end for overlap.
  \item \textbf{Search ceiling.} The top-cited ranking only considers
        papers in the Semantic Scholar keyword-search top 1{,}000 per
        query, not the full literature.
  \item \textbf{No field normalisation.} We use raw \texttt{citationCount},
        not OpenAlex's field-weighted citation impact (FWCI). Raw counts
        overweight older venues; the two-pass benchmark partially
        mitigates this but does not eliminate it.
  \item \textbf{Widened review labels.} Accepting \emph{survey} and
        \emph{literature review} alongside \emph{systematic \ldots} raises
        recall but admits narrative reviews without protocol documentation.
        A study scoped to SLR-quality would need a stricter gate.
  \item \textbf{Manual classification subjectivity.} The SLR-vs.-narrative-survey
        boundary is fuzzy; another reviewer may classify differently. Our
        rubric and per-paper justifications are in
        \texttt{data/manual/ss/slr\_decisions.csv}.
  \item \textbf{Subfield boundary.} Choosing ``technical debt'' as the unit
        aggregates over sub-strands (self-admitted technical debt,
        architectural debt, requirements technical debt) that may behave
        differently.
  \item \textbf{English-language indexing and metadata gaps.} Non-English
        studies were excluded; empty-bibliography SLRs were excluded by
        construction (50 of 110 auto-INCLUDEs), which biases the corpus
        toward better-indexed venues.
  \item \textbf{Citation counts are point-in-time.} The date control
        mitigates the eligibility filter but not the underlying ranking
        volatility.
  \item \textbf{Coverage $\neq$ validity.} Low coverage does not prove
        incorrect SLR conclusions.
\end{itemize}

\section{Conclusion}

Within the technical-debt subfield, published SLRs cite, on average,
10.5\,\% of the date-eligible top-50 most-cited papers, with a median of
4.4\,\%. Forty percent of SLRs cite none of the eligible top-50. Gaps are
systematic---concentrated in journal venues (90.4\,\% of misses), peaking
at both 6--10-year-old foundations and 0--2-year-old recent work, and not
primarily explained by paywalls (43\,\% of missed papers are open-access).
The findings are robust under a top-100 benchmark (mean drops only to
8.6\,\%). A small set of protocol additions---backward snowballing from
canonical seeds, citation-aware cross-checks against a top-$N$ benchmark,
multi-source deduplication, and explicit time-window framing---would
materially close the gap. Readers of SLRs in this subfield should not
assume comprehensive coverage by default; reviewers and editors should
ask SLR authors to demonstrate it.

\section*{Acknowledgments}

The authors thank Prof.\ Ian Gorton for advising the project and for
methodological feedback that shaped the two-pass top-cited corpus design
(\S\ref{sec:coverage}, \S2.5). We also acknowledge the Semantic Scholar
Open API and the Crossref REST API for free access to the citation and
reference-graph data this study depends on.

\bibliographystyle{IEEEtran}
\bibliography{references}

\section*{Replicability}

All code, classifications, and intermediate data live in the project
repository. The full SS-only pipeline can be re-run end-to-end via
\texttt{PYTHONPATH=. python sources/ss/run.py} after installing
dependencies and setting \texttt{SEMANTIC\_SCHOLAR\_API\_KEY} and
\texttt{CROSSREF\_MAILTO}. Static figures are regenerated via
\texttt{python report/build\_figures.py --source ss}. The interactive
overlap explorer is launched with \texttt{python tools/explorer/serve.py}.

\end{document}
